%% (H1) Základy umělé inteligence
\subsection{Řešení úloh prohledáváním}
\begin{itemize}
\item formulace úlohy
\begin{itemize}
\item agent - vnímá prostředí pomocí senzorů, aktuátory mu umožňují jednat
\begin{itemize}
\item racionální agent
\begin{itemize}
\item zvolí akci, která maximalizuje jeho výkon
\end{itemize}
\item simple reflex agent
\begin{itemize}
\item observation $\to$ action
\end{itemize}
\item model-based reflex agent
\begin{itemize}
\item past state + past action + observation $\to$ state
\item state $\to$ action
\end{itemize}
\item goal-based agent
\begin{itemize}
\item state + goal $\to$ action
\end{itemize}
\end{itemize}
\item reprezentace stavů agenta
\begin{itemize}
\item atomic - stav je blackbox
\begin{itemize}
\item dá se použít prohledávání
\end{itemize}
\item factored - stav je vektor
\begin{itemize}
\item dá se použít constraint satisfaction, výroková logika, plánování
\end{itemize}
\item structured - stav je sada objektů, ty jsou propojeny
\begin{itemize}
\item dá se použít (predikátová) logika prvního řádu
\end{itemize}
\end{itemize}
\item problem solving agent - typ goal-based agenta
\begin{itemize}
\item atomická reprezentace stavů
\item cíl = množina cílových stavů
\item akce = přechody mezi stavy
\item hledáme posloupnost akcí, kterými se dostaneme z výchozího do nějakého cílového stavu
\item očekáváme, že prostředí je plně pozorovatelné, diskrétní, statické a deterministické
\item příklad - Lloydova patnáctka
\begin{itemize}
\item ne všechny stavy jsou dosažitelné - zachovává se parita permutace
\end{itemize}
\end{itemize}
\item formulace problému
\begin{itemize}
\item budeme potřebovat abstrakci prostředí (odstraníme detaily z prostředí)
\begin{itemize}
\item validní abstrakce = abstraktní řešení můžeme expandovat do reálného řešení
\item užitečná abstrakce = vykonávání akcí v řešení je jednodušší než v původním problému
\end{itemize}
\item dobře definovaný problém
\begin{itemize}
\item výchozí stav
\item přechodový model \texttt{(\allowbreak{}state,\allowbreak{} \allowbreak{}action)\allowbreak{} \allowbreak{}-\textgreater{}\allowbreak{} \allowbreak{}state}
\item goal test (funkce, která odpoví true/false podle toho, zda je daný stav cílový)
\end{itemize}
\end{itemize}
\end{itemize}
\item stromové vs. grafové prohledávání
\begin{itemize}
\item udržujeme si hranici (frontier) prozkoumané oblasti
\item graph search si navíc ukládá \quotedblbase{}uzavřenost\textquotedblleft{} vrcholu (jestli už jsme vrchol viděli)
\begin{itemize}
\item takže nenastávají problémy s cykly
\item kvůli tomu, že stavový prostor může být nekonečný, neinicializujeme všechny vrcholy, ale místo toho použijeme hešovací tabulku
\end{itemize}
\item i tree search lze použít k prohledávání grafů s cykly, ale může se chovat divně
\item u prohledávání DAGů tree search nezaručuje, že se prohledané stavy nebudou opakovat (opakovat se můžou, pokud do stejného vrcholu vedou dvě orientované cesty)
\end{itemize}
\item neinformované prohledávání (DFS, BFS, uniform-cost search)
\begin{itemize}
\item BFS (fronta)
\begin{itemize}
\item nejmělčí neexpandovaný vrchol je zvolen k expanzi
\item úplný pro konečně větvící grafy
\item optimální - pokud je cena cesty nerostoucí funkce hloubky
\item časová a prostorová složitost $O(b^{d+1})$, kde $b$ je stupeň větvení (branching factor, maximální výstupní stupeň) a $d$ je hloubka nejbližšího cíle
\item problém je s prostorovou složitostí
\end{itemize}
\item DFS (zásobník)
\begin{itemize}
\item nejhlubší neexpandovaný vrchol je zvolen k expanzi
\item úplný pro graph search
\begin{itemize}
\item úplný = nalezne řešení vždy, když existuje (a v opačném případě hlásí, že řešení neexistuje)
\end{itemize}
\item neúplný pro tree search (pokud algoritmus pustíme na grafu, který není strom)
\item suboptimální - nesměřuje k cíli
\item časová složitost $O(b^m)$, kde $m$ je maximální hloubka libovolného vrcholu
\item prostorová složitost $O(bm)$
\item pokud použijeme backtracking (generujeme jenom jednoho následníka místo všech), tak nám stačí dokonce jenom $O(m)$ paměti
\begin{itemize}
\item DFS využívá to, že nám stav může vrátit všechny následníky (pak je držíme v~paměti)
\item backtracking v daném vrcholu generuje jednoho následníka (takže všechny ostatní nemusíme držet v paměti)
\begin{itemize}
\item ale někdy se nedá generovat jenom jeden následník
\end{itemize}
\end{itemize}
\end{itemize}
\item rozšíření BFS o funkci určující cenu kroku (tedy cenu hrany)
\begin{itemize}
\item Dijkstrův algoritmus
\item taky se tomu říká uniform-cost search
\item $g(n)$ označuje cenu nejlevnější cesty ze startu do $n$
\item místo fronty použijeme prioritní frontu
\end{itemize}
\end{itemize}
\item informované prohledávání (algoritmus A*, přípustné a konzistentní heuristiky)
\begin{itemize}
\item best-first search
\begin{itemize}
\item pro vrchol máme ohodnocovací funkci $f(n)$
\item $f(n)$ použijeme v Dijkstrově algoritmu místo $g(n)$
\item máme heuristickou funkci $h(n)$, která pro daný vrchol označuje odhadovanou cenu nejlevnější cesty do cílového stavu
\item best-first je třída algoritmů, kde máme uzly ohodnoceny nějakou funkcí a vybíráme nejmenší z nich
\item greedy best-first search
\begin{itemize}
\item $f(n)=h(n)$
\item není optimální ani úplný
\end{itemize}
\item A* algoritmus
\begin{itemize}
\item $f(n)=g(n)+h(n)$, kde $h(n)$ je heuristika
\item optimální a úplný
\item typicky mu dojde paměť dřív než čas
\end{itemize}
\end{itemize}
\item co chceme od heuristiky $h(n)$ u algoritmu A*
\begin{itemize}
\item přípustnost - $h(n)$ musí být menší rovna nejkratší cestě z daného vrcholu do cíle, musí být nezáporná
\item monotónnost (= konzistence)
\begin{itemize}
\item mějme vrchol $n$ a jeho souseda $n'$
\item $c(n,a,n')$ je cena přechodu z $n$ do $n'$ (akcí $a$)
\item heuristika je monotónní, když $h(n)\leq c(n,a,n')+h(n')$
\end{itemize}
\item tvrzení: monotonie implikuje přípustnost, opačná implikace neplatí
\item tvrzení: pro monotónní heuristiku jsou hodnoty $f(n)$ neklesající po libovolné cestě
\item tvrzení: pokud je $h(n)$ přípustná, pak je A* u tree search optimální
\begin{itemize}
\item do cíle nevedou žádné suboptimální cesty, takže si vystačíme s důkazem optimality Dijkstrova algoritmu
\end{itemize}
\item tvrzení: pokud je $h(n)$ monotónní, pak je A* u graph search optimální
\begin{itemize}
\item s nemonotónní heuristikou jsme mohli najít suboptimální cestu do cíle
\end{itemize}
\end{itemize}
\item když heuristika $h_2$ dává větší hodnoty než $h_1$, tak se říká, že $h_2$ dominuje $h_1$
\begin{itemize}
\item pokud je $h_2$ přípustná a pokud se $h_2$ nepočítá výrazně déle než $h_1$, tak je $h_2$ zjevně lepší než $h_1$
\end{itemize}
\end{itemize}
\end{itemize}
\subsection{Splňování omezujících podmínek}
\begin{itemize}
\item problém splňování podmínek
\begin{itemize}
\item = constraint satisfaction problem (CSP)
\item konečná množina proměnných
\item domény - konečné množiny možných hodnot pro každou proměnnou
\item konečná množina podmínek (constraints)
\begin{itemize}
\item constraint je relace na podmnožině proměnných (je to podmnožina kartézského součinu domén)
\item constraint arity = počet proměnných, které podmínka omezuje
\end{itemize}
\item chceme přípustné řešení (feasible solution)
\begin{itemize}
\item úplné konzistentní přiřazení hodnot proměnným
\item konzistentní \dots{} všechny podmínky jsou splněny
\end{itemize}
\item CSP můžeme řešit tree-search backtrackingem
\begin{itemize}
\item úpravou proměnných v určitém pořadí můžeme získat větší efektivitu
\item můžeme pročistit hodnoty, které zjevně nesplňují podmínky
\begin{itemize}
\item příklad
\begin{itemize}
\item $a\in\set{3,4,5,6,7}$
\item $b\in\set{1,2,3,4,5}$
\item constraint: $a\lt b$
\item ale např. $a=6$ vůbec nedává smysl
\item odstraněním nekonzistentních hodnot dostaneme $a\in \set{3,4}$, $b\in\set{4,5}$
\end{itemize}
\item můžeme použít metodu hranové konzistence
\end{itemize}
\end{itemize}
\end{itemize}
\item hranová konzistence (algoritmus AC-3)
\begin{itemize}
\item = arc consistency (arc \dots{} orientovaná hrana)
\item uvažujme binární podmínky
\begin{itemize}
\item libovolnou n-ární podmínku lze převést na balík binárních podmínek
\end{itemize}
\item každá binární podmínka odpovídá dvojici orientovaných hran v síti podmínek (vrcholy odpovídají proměnným)
\begin{itemize}
\item pokud je mezi dvěma proměnnými více podmínek, můžeme je sloučit do jedné (abychom neměli multigraf)
\end{itemize}
\item orientovaná hrana $(V_i,V_j)$ je hranově konzistentní $\equiv$ pro každé $x\in D_i$ existuje $y\in D_j$ takové, že přiřazení $V_i=x$ a $V_j=y$ splní všechny binární podmínky pro $V_i,V_j$
\begin{itemize}
\item může platit, že $(V_i,V_j)$ je hranově konzistentní, ale $(V_j,V_i)$ není
\end{itemize}
\item CSP je hranově konzistentní $\equiv$ každá hrana je hranově konzistentní (v obou směrech)
\item algoritmus AC-3
\begin{itemize}
\item začínám se všemi hranami (podmínkami) ve frontě
\item postupně beru hrany z fronty a pro každou odstraním nekonzistentní prvky z~domény
\item pokud jsem odstranil nějaké prvky z domény, musím zopakovat kontrolu pro hrany směřující do dané proměnné (kromě hrany opačné k té aktuálně zpracované)
\item složitost $O(cd^3)$, kde $c$ je počet podmínek, $d$ je velikost domény
\begin{itemize}
\item $O(d^2)$ \dots{} kontrola konzistence
\item každá podmínka se ve frontě může objevit nejvýše $d$-krát, protože se z dané domény dá odstranit nejvýše $d$ prvků
\end{itemize}
\item optimální algoritmus má složitost $O(cd^2)$
\end{itemize}
\item hranová konzistence je typ lokální konzistence, negarantuje globální konzistenci
\item arc consistency (AC) vs. forward checking (FC)
\begin{itemize}
\item FC je jednodušší (a rychlejší) než AC, slouží ke stejnému účelu - zrychlení prohledávání
\item FC na základě přiřazení do proměnné (v daném kroku prohledávání) pročistí domény dosud nepřiřazených proměnných
\item AC k tomu navíc zajišťuje vzájemnou kompatibilitu domén nepřiřazených proměnných
\end{itemize}
\end{itemize}
\item algoritmus hledání řešení (MAC)
\begin{itemize}
\item chceme zkombinovat AC a backtracking
\item problém uděláme hranově konzistentní
\item po každém přiřazení obnovíme hranovou konzistenci
\item této technice se říká \textit{look ahead} nebo \textit{constraint propagation} nebo \textit{udržování hranové konzistence} (MAC, maintaining arc consistency)
\item kontroly FC/AC jsou v polynomiálním čase, kdežto strom stavů se větví exponenciálně
\end{itemize}
\end{itemize}
\subsection{Logické uvažování}
\begin{itemize}
\item základy výrokové logiky (konjunktivní a disjunktivní normální forma)
\begin{itemize}
\item viz \href{spolecna-matematika.md}{Společná matematika}
\item výrokovou logiku lze použít k logickému odvozování znalostí na základě znalostní báze
\end{itemize}
\item algoritmus DPLL
\begin{itemize}
\item literál $\ell$ má čistý výskyt ve $\varphi$, pokud se $\ell$ vyskytuje ve $\varphi$ a opačný literál $\overline\ell$ se ve $\varphi$ nevyskytuje
\begin{itemize}
\item takový literál můžu nastavit na 1
\item to neovlivní splnitelnost výroku, ale zmenší to množinu modelů, které jsem schopen nalézt
\end{itemize}
\item algoritmus
\begin{itemize}
\item dokud $\varphi$ obsahuje jednotkovou klauzuli $\ell$, ohodnoť $\ell=1$ a proveď jednotkovou propagaci
\begin{itemize}
\item každou klauzuli obsahující $\ell$ odstraníme (protože je takto splněna)
\item $\overline\ell$ odstraníme ze všech klauzulí, které ho obsahují (protože $\overline\ell$ nemůže zajistit splnění dané klauzule)
\end{itemize}
\item dokud existuje literál $\ell$, který má ve $\varphi$ čistý výskyt, ohodnoť $\ell=1$ a odstraň klauzule obsahující $\ell$
\item pokud $\varphi$ neobsahuje žádnou klauzuli, je splnitelný
\item pokud $\varphi$ obsahuje prázdnou klauzuli, není splnitelný
\item jinak zvol dosud neohodnocenou výrokovou proměnnou $p$ a zavolej algoritmus rekurzivně na $\varphi\land p$ a na $\varphi\land\neg p$
\end{itemize}
\item algoritmus běží v exponenciálním čase
\item praktické poznámky
\begin{itemize}
\item čistý výskyt se zas tak moc nepoužívá
\item jednotková propagace
\begin{itemize}
\item hledám klauzule o jedné proměnné
\item je v zásadě ekvivalentní hranové konzistenci
\end{itemize}
\item ke zkoušení hodnot proměnných se používá tree search
\end{itemize}
\end{itemize}
\item dopředné a zpětné řetězení (Hornovské klauzule)
\begin{itemize}
\item mám dotaz $\alpha$, ten vyjádřím jako výrokovou formuli
\begin{itemize}
\item znalostní bázi $KB$ vyjádřím jako teorii
\item zajímá mě, zda $KB\models\alpha$ (to platí, právě když $KB\land\neg\alpha$ je nesplnitelné)
\end{itemize}
\item řetězení můžeme použít, pokud znalostní báze obsahuje pouze Hornovy klauzule
\begin{itemize}
\item Hornova klauzule = disjunkce literálů, z nichž nejvýše jeden je pozitivní
\item Hornovu klauzuli $\neg p\lor\neg q\lor\neg r\lor s$ lze zapsat jako implikaci $p\land q\land r\implies s$
\end{itemize}
\item forward chaining \dots{} data-driven reasoning
\begin{itemize}
\item $p\land q\land r\implies s$
\item pokud vím, že platí $p$, pak převedu na $q\land r\implies s$
\item jakmile se počet předpokladů sníží na 0, vím, že $s$ platí
\item je to v podstatě speciální případ použití rezolučního pravidla
\begin{itemize}
\item rezolvuju to, co vím, s libovolnou klauzulí
\end{itemize}
\item algoritmus
\begin{itemize}
\item každá proměnná má počítadlo proměnných, které na ni ukazují
\item postupně beru pravdivé proměnné, snižuju počítadla u proměnných, na které ukazují
\item jakmile u proměnné klesne počítadlo na nulu, zařadím mezi pravdivé proměnné
\end{itemize}
\end{itemize}
\item backward chaining \dots{} goal-driven reasoning
\begin{itemize}
\item něco mě zajímá - pokouším se to odvodit
\begin{itemize}
\item najdu klauzule, kde je nějaký literál z $\alpha$ na pravé straně implikace
\item tak dostanu další literály, které musím dokázat
\item postupuju tak dlouho, dokud se nedostanu k faktům (literálům, o nichž vím, že platí)
\end{itemize}
\item tohle se používá v Prologu
\item opět vlastně používám rezoluční pravidlo
\begin{itemize}
\item to, co chci zjistit, rezolvuju s libovolnou klauzulí
\end{itemize}
\end{itemize}
\item forward chaining prochází spoustu klauzulí, které nás nezajímají - backward chaining může mít výrazně menší složitost
\end{itemize}
\item rezoluce
\begin{itemize}
\item mám dotaz $\alpha$, ten vyjádřím jako výrokovou formuli
\begin{itemize}
\item znalostní bázi $KB$ vyjádřím jako teorii
\item zajímá mě, zda $KB\models\alpha$ (to platí, právě když $KB\land\neg\alpha$ je nesplnitelné)
\end{itemize}
\item $KB\land\neg\alpha$ převedeme do CNF $\to$ dostaneme množinu klauzulí (formuli $S$)
\item použijeme rezoluci (zkonstruujeme rezoluční zamítnutí formule $S$)
\begin{itemize}
\item použitím rezolučního pravidla z klauzulí $\varphi\lor p$ a $\psi\lor\neg p$ dostaneme $\varphi\lor\psi$
\item zkoušíme spolu rezolvovat postupně všechny dvojice klauzulí
\item pokud v průběhu dostaneme prázdnou klauzuli, máme rezoluční důkaz $\alpha$ z $KB$
\item pokud se dostaneme do stavu, kdy nelze použít rezoluční pravidlo na žádnou dvojici klauzulí, tak víme, že neplatí $KB\models\alpha$
\end{itemize}
\end{itemize}
\end{itemize}
\subsection{Pravděpodobnostní uvažování}
\begin{itemize}
\item základy teorie pravděpodobnosti (úplná sdružená distribuce, nezávislost, Bayesovo pravidlo)
\begin{itemize}
\item viz \href{spolecna-matematika.md}{Společná matematika}
\item pravděpodobnostní uvažování umožňuje agentovi zacházet s nejistotou
\item zdroji nejistoty jsou částečná pozorovatelnost (nelze snadno zjistit současný stav) a nedeterminismus (není jisté, jak akce dopadne)
\item agent si pamatuje belief state - pravděpodobnostní distribuce přes všechny možné stavy světa, v nichž se agent může nacházet
\end{itemize}
\item pravděpodobnostní uvažování (vysčítání, normalizace)
\begin{itemize}
\item někdy si můžeme udělat tabulku všech možností a určit jejich pravděpodobnosti (full joint probability distribution)
\begin{itemize}
\item odpověď lze vyjádřit z ní (této metodě se říká enumeration, výčet) - ale u velkých tabulek je někdy potřeba posčítat příliš mnoho hodnot
\end{itemize}
\item $P(Y)=\sum_z P(Y\mid z)\cdot P(z)$
\item $P(Y\mid e)=\frac{P(Y\land e)}{P(e)}$
\begin{itemize}
\item typicky nemusíme znát $P(e)$, stačí použít $\alpha P(Y\land e)$, kde $\alpha$ je normalizační konstanta
\end{itemize}
\item velkou tabulku můžeme někdy reprezentovat menším způsobem - pokud jsou proměnné podmíněně nezávislé
\item Bayesovo pravidlo \dots{} $P(Y\mid X)=\frac{P(X\mid Y)P(Y)}{P(X)}=\alpha P(X\mid Y) P(Y)$
\begin{itemize}
\item lze odvodit z definice podmíněné pravděpodobnosti
\item pomůže nám při přechodu z příčinného k diagnostickému směru
\begin{itemize}
\item $P(\mathrm{effect} \mid \mathrm{cause})$ \dots{} casual direction (příčinný směr), obvykle známý
\item $P(\mathrm{cause} \mid \mathrm{effect})$ \dots{} diagnostic direction (diagnostický směr), obvykle chceme zjistit
\item známe důsledky; chceme zjistit, jaké příčiny k nim mohly vést
\end{itemize}
\end{itemize}
\item naivní bayesovský model
\begin{itemize}
\item $P(\text{Cause}, \text{Effect}_1, \dots, \text{Effect}_n) = P(\text{Cause})\prod_iP(\text{Effect}_i\mid\text{Cause})$
\item předpokládáme nezávislost důsledků
\end{itemize}
\item chain rule
\begin{itemize}
\item $P(x_1,\dots,x_n)=\prod_iP(x_i\mid x_{i-1},\dots,x_1)$
\item využívá product rule (lze odvodit z definice podmíněné pravděpodobnosti)
\end{itemize}
\end{itemize}
\item Bayesovské sítě: konstrukce a vztah k úplné sdružené distribuci
\begin{itemize}
\item bayesovská síť je DAG, kde vrcholy odpovídají náhodným proměnným
\begin{itemize}
\item šipky popisují závislost
\item u každého vrcholu jsou CPD tabulky, které popisujou jeho závislost na rodičích
\begin{itemize}
\item CPD = conditional probability distribution
\item pro každou kombinaci hodnot rodičů jsou v CPD tabulce pravděpodobnosti hodnot dané náhodné proměnné
\end{itemize}
\end{itemize}
\item \begin{center}\includegraphics[max width=0.85\linewidth,max height=8cm]{notes-ipp/statnice/prilohy/bayesovska-sit.png}\end{center}
\begin{itemize}
\item autorem obrázku je \href{https://ktiml.mff.cuni.cz/~bartak/ui_intro/}{Roman Barták}
\item proměnné jsou binární, proto má tabulka vždy jen jeden sloupec (druhý sloupec by mohl obsahovat pravděpodobnost, že jev nenastane - lze ho dopočítat z prvního sloupce)
\end{itemize}
\item konstrukce bayesovské sítě
\begin{itemize}
\item nějak si uspořádáme proměnné
\begin{itemize}
\item lepší je řadit je od příčin k důsledkům - budeme mít menší síť a CPD tabulky se nám budou vyplňovat snáze
\end{itemize}
\item jdeme odshora dolů, přidáváme správné hrany
\item příklad: MaryCalls, JohnCalls, Alarm, Burglary, Earthquake (záměrně jdeme od důsledků k příčinám, byť to není praktické)
\begin{itemize}
\item z MaryCalls povede hrana do JohnCalls, protože nejsou nezávislé (když volá Mary, tak je větší šance, že volá i John)
\item z obou povede hrana do Alarmu
\item z Alarmu vedou hrany do Burglary a Earthquake (ale hrany z JohnCalls a MarryCalls tam nepovedou - na těch hranách nezáleží, jsou nezávislé)
\item z Burglary povede hrana do Earthquake, protože pokud zní alarm a k vloupání nedošlo, tak pravděpodobně došlo k zemětřesení
\item akorát se nám v tomto pořadí budou špatně určovat CPD
\end{itemize}
\end{itemize}
\item bayesovská síť zachycuje úplnou sdruženou (joint) distribuci
\begin{itemize}
\item dokážeme určit pravděpodobnost libovolné kombinace hodnot náhodných proměnných
\item pronásobíme mezi sebou hodnoty z tabulek
\end{itemize}
\end{itemize}
\item Bayesovské sítě: exaktní odvozování (enumerace, eliminace proměnných)
\begin{itemize}
\item z bayesovských sítí můžeme provádět inferenci - odvozovat pravděpodobnost proměnných pomocí pravděpodobností \textit{skrytých} proměnných
\item $P(X\mid e)=\alpha P(X,e)=\alpha\sum_y P(X,e,y)$
\begin{itemize}
\item kde pravděpodobnost $X$ odvozujeme, $e$ jsou pozorované proměnné (evidence) a $y$ jsou hodnoty další skryté proměnné $Y$
\end{itemize}
\item přičemž $P(X,e,y)$ lze určit pomocí CPD tabulek
\begin{itemize}
\item když nemůžeme určit konkrétní hodnotu pravděpodobnosti přímo, tak výpočet rozvětvíme
\item některé větve se objeví vícekrát - hodí se nám dynamické programování
\end{itemize}
\item používáme \textit{faktory} (tabulky odpovídající CPD tabulkám)
\begin{itemize}
\item tabulky s faktory mezi sebou můžeme násobit
\item eliminace
\begin{itemize}
\item mějme faktor s pravděpodobnostmi $P(A,B,C)$ pro všechny možné kombinace hodnot proměnných
\item můžeme ho rozdělit na dvě půlky (v jedné bude všude platit $A$, v druhé $\neg A$)
\item tyhle dvě půlky můžeme sečíst (vždy odpovídající řádky)
\item tak eliminujeme proměnnou $A$, dostaneme faktor s pravděpodobnostmi $P(B,C)$
\end{itemize}
\end{itemize}
\end{itemize}
\item Bayesovské sítě: aproximační odvozování (Monte Carlo, zamítání, vážení věrohodností)
\begin{itemize}
\item odhadujeme hodnoty, které je těžké přímo spočítat
\begin{itemize}
\item třeba když je síť moc velká a hustě propojená
\item zajímá nás $P(X\mid e)$
\end{itemize}
\item generujeme hodně vzorků, použijeme statistiku
\item pravděpodobnostní rozdělení náhodných veličin známe (postupujeme topologicky bayesovskou sítí a generujeme hodnoty)
\item jak generovat vzorky?
\begin{itemize}
\item rejection sampling = zahodíme vzorky, kde neplatí $e$
\begin{itemize}
\item odhadneme $P(X\mid e)\approx\frac{N(X,e)}{N(e)}$
\begin{itemize}
\item $N(X,e)$ \dots{} počet vzorků, kde platí $X,e$
\end{itemize}
\item riziko, že zahodíme příliš mnoho vzorků
\end{itemize}
\item likelihood weighting
\begin{itemize}
\item zafixujeme (ignorujeme) $e$, generujeme jen zbylé proměnné
\item počtům $N$ přidělíme váhy podle pravděpodobnosti $e$ (např. místo toho, abychom zahodili cca 90 \% vzorků, přiřadíme těm počtům váhu 0.1)
\item problém nastává, když jsou váhy příliš malé
\end{itemize}
\end{itemize}
\end{itemize}
\end{itemize}
\subsection{Reprezentace znalostí}
\begin{itemize}
\item situační kalkulus, problém rámce
\begin{itemize}
\item problém rámce
\begin{itemize}
\item návrh: stav světa/agenta budeme popisovat pomocí časově anotovaných výrokových proměnných
\item musíme definovat spoustu axiomů, které zajistí, že se hodnoty proměnných mezi časovými kroky nemění náhodně
\item tak bychom snadno zahltili odvození pomocí výrokové logiky, protože bychom měli hodně akcí a stavů
\end{itemize}
\item jak předejít opakování axiomů pro různé časy a podobné akce?
\begin{itemize}
\item použijeme predikátovou logiku prvního řádu (\textit{situation calculus})
\item náš modelovaný svět se vyvíjí - postupně nastávají různé \textit{situace}
\begin{itemize}
\item úvodní situace \dots{} $s_0$
\item situace po aplikace akce $a$ na situaci $s$ \dots{} $\text{Results}(s,a)$
\end{itemize}
\item stavy jsou definovány jako relace mezi objekty
\begin{itemize}
\item pokud se relace mění, musíme přidat další argument označující situaci, v níž relace platí: $\mathrm{at}(\mathrm{robot}, \mathrm{location}, s)$
\item pro stálé relace takový argument není potřeba: $\mathrm{connected}(l, l')$
\end{itemize}
\item podmínky (předpoklady) akcí jsou definovány possibility axiomem
\begin{itemize}
\item obecný tvar: $\varphi(s)\implies \text{Poss}(a,s)$
\begin{itemize}
\item kde $\varphi$ je formule popisující předpoklady akce $a$
\end{itemize}
\item např. $\mathrm{at}(a, l, s) \land \mathrm{connected}(l, l') \implies \mathrm{Poss}(\mathrm{go}(a, l, l'), s)$
\end{itemize}
\item vlastnosti dalšího stavu jsou popsány pomocí successor-state axiomů
\begin{itemize}
\item např. $\mathrm{Poss}(a, s) \implies$ $(\mathrm{at}(a,y,\mathrm{Results}(s, a)) \iff a=\mathrm{go}(a,x,y) \lor (\mathrm{at}(a,y,s) \land a\neq \mathrm{go}(a,y,z))$
\end{itemize}
\item plánování se provádí tak, že se zeptáme, zda existuje situace, kde je cílová podmínka pravdivá
\end{itemize}
\end{itemize}
\item Markovské modely - filtrace, predikce, vyhlazování, nejpravděpodobnější průchod
\begin{itemize}
\item každý stav světa je popsán množinou náhodných proměnných
\begin{itemize}
\item skryté náhodné proměnné $X_t$ - popisují reálný stav
\item pozorovatelné náhodné proměnné $e_t$ - popisují, co pozorujeme
\item uvažujeme diskrétní čas, takže $t$ označuje konkrétní okamžik
\item množinu proměnných od $X_a$ do $X_b$ budeme značit jako $X_{a:b}$
\end{itemize}
\item formální model
\begin{itemize}
\item skrytý Markovův model (HMM) je dvojice $(X_n,e_n)$
\begin{itemize}
\item $X_n$ a $e_n$ jsou náhodné (stochastické) procesy, takže vlastně posloupnosti náhodných proměnných
\item $X_n$ je markovský proces, který nemůžeme přímo pozorovat
\end{itemize}
\item transition model
\begin{itemize}
\item určuje pravděpodobnostní rozložení proměnných posledního stavu, když známe jejich předchozí hodnoty \dots{} $P(X_t\mid X_{0:t-1})$
\item zjednodušující předpoklady
\begin{itemize}
\item další stav závisí jen na předchozím stavu \dots{} \textit{Markov assumption}
\item všechny přechodové tabulky $P(X_t\mid X_{t-1})$ jsou identické přes všechna $t$ \dots{} \textit{stationary process}
\end{itemize}
\end{itemize}
\item sensor (observation) model
\begin{itemize}
\item popisuje, jak pozorované proměnné závisí na ostatních proměnných \dots{} $P(e_t\mid X_{0:t},e_{1:t-1})$
\item zjednodušující předpoklad: pozorování záleží jen na aktuálním stavu \dots{} \textit{sensor Markov assumption}
\end{itemize}
\end{itemize}
\item základní inferenční úlohy
\begin{itemize}
\item filtrování - znám minulá pozorování, zajímá mě přítomný stav
\begin{itemize}
\item $P(X_{t+1}\mid e_{1:t+1})=P(X_{t+1}\mid e_{1:t},e_{t+1})$
\begin{itemize}
\item použijeme Bayesovo pravidlo
\end{itemize}
\item $=\alpha\cdot P(e_{t+1}\mid X_{t+1},e_{1:t})\cdot P(X_{t+1}\mid e_{1:t})$
\begin{itemize}
\item použijeme \textit{sensor Markov assumption}
\end{itemize}
\item $=\alpha\cdot P(e_{t+1}\mid X_{t+1})\cdot P(X_{t+1}\mid e_{1:t})$
\item $=\alpha\cdot P(e_{t+1}\mid X_{t+1})\cdot \sum_{x_t}P(X_{t+1}\mid x_t,e_{1:t})\cdot P(x_t\mid e_{1:t})$
\item $=\alpha\cdot P(e_{t+1}\mid X_{t+1})\cdot \sum_{x_t}P(X_{t+1}\mid x_t)\cdot P(x_t\mid e_{1:t})$
\item užitečný filtrovací algoritmus si musí udržovat odhad současného stavu, jelikož je vzorec rekurzivní
\item používá se \textit{forward} algoritmus
\end{itemize}
\item předpověď (prediction) - znám minulá pozorování, zajímají mě budoucí stavy
\begin{itemize}
\item vlastně jako filtering, akorát nepřidávám další pozorování (měření)
\item $P(X_{t+k+1}\mid e_{1:t})=\sum_{x_{t+k}} P(X_{t+k+1}\mid x_{t+k})\cdot P(x_{t+k}\mid e_{1:t})$
\item po určitém čase (\textit{mixing time}) distribuce předpovědí konverguje ke stacionárnímu rozdělení daného markovského procesu a zůstane konstantní
\end{itemize}
\item vyhlazování (smoothing) - znám minulá pozorování, zajímají mě minulé stavy
\begin{itemize}
\item $P(X_k\mid e_{1:t})=P(X_k\mid e_{1:k},e_{k+1:t})$
\begin{itemize}
\item použijeme Bayesovo pravidlo
\end{itemize}
\item $=\alpha\cdot P(X_k\mid e_{1:k})\cdot P(e_{k+1:t}\mid X_k,e_{1:k})$
\begin{itemize}
\item použijeme \textit{sensor Markov assumption}
\end{itemize}
\item $=\alpha\cdot P(X_k\mid e_{1:k})\cdot P(e_{k+1:t}\mid X_k)$
\item přitom levý člen známe z filteringu (je to \quotedblbase{}dopředný směr\textquotedblleft{}), pravý je \quotedblbase{}zpětný směr\textquotedblleft{} z naměřených hodnot \textit{v budoucnosti} (relativně vůči danému okamžiku)
\item používá se \textit{forward-backward} algoritmus
\end{itemize}
\item nejpravděpodobnější vysvětlení (most likely explanation) - znám minulá pozorování, zajímá mě nejpravděpodobnější posloupnost minulých stavů
\begin{itemize}
\item hledáme posloupnost $X_{1:t}$ takovou, že má největší $P(X_{1:t+1}\mid e_{1:t+1})$
\item opět existuje rekurzivní vzorec - podíváme se na pravděpodobnosti předchozích stavů a na nejpravděpodobnější \quotedblbase{}cesty\textquotedblleft{} do těchto stavů
\item používá se \textit{Viterbiho} algoritmus
\end{itemize}
\end{itemize}
\end{itemize}
\item skryté Markovské modely (HMM) vs. dynamické Bayesovské sítě
\begin{itemize}
\item skryté Markovovy modely
\begin{itemize}
\item předpokládejme, že stav procesu je popsán jedinou diskrétní náhodnou proměnnou $X_t$ (a máme jednu proměnnou $E_t$ odpovídající pozorování)
\item pak můžeme všechny základní algoritmy (filtering, prediction, smoothing, \dots{}) implementovat maticově
\end{itemize}
\item dynamická bayesovská síť
\begin{itemize}
\item umožňuje popsat víc náhodných proměnných
\item reprezentuje časový pravděpodobnostní model
\item zachycuje vztah mezi minulým a současným časovým okamžikem (minulou a současnou vrstvou)
\item každá stavová proměnná má rodiče ve stejné vrstvě nebo v předchozí (podle Markovova předpokladu)
\end{itemize}
\item dynamická bayesovská síť (DBN) vs. skrytý Markovův model (HMM)
\begin{itemize}
\item skrytý Markovův model je speciální případ dynamické bayesovské sítě
\item dynamická bayesovská síť může být kódovaná jako skrytý Markovův model
\begin{itemize}
\item jedna náhodná proměnná ve skrytém Markovově modelu je n-tice hodnot stavových proměnných v dynamické bayesovské síti
\item v HMM vlastně celý stav světa komprimujeme do jedné náhodné proměnné - v DBN jich můžeme mít víc
\end{itemize}
\item vztah mezi DBN a HMM je podobný jako vztah mezi běžnými bayesovskými sítěmi a tabulkou s \textit{full joint probability distribution}
\begin{itemize}
\item DBN je úspornější
\end{itemize}
\end{itemize}
\end{itemize}
\end{itemize}
\subsection{Automatické plánování}
\begin{itemize}
\item formulace plánovacího problému (definice operátoru)
\begin{itemize}
\item stav světa je reprezentován jako množina proměnných
\begin{itemize}
\item stav je popsán atomy, které platí (ostatní neplatí, používáme předpoklad uzavřeného světa)
\item pravdivnostní hodnota některých atomů se mění \dots{} fluents
\item u jiných atomů je pravdivostní hodnota stálá \dots{} rigid atoms
\end{itemize}
\item schémata akcí (operátory) popisují možnosti agenta
\begin{itemize}
\item schéma akce (operátor) popisuje akci, aniž by specifikovalo objekty, na které se akce vztahuje
\item operátor je trojice (název, předpoklady, důsledky)
\begin{itemize}
\item název - obsahuje rovněž parametry operátoru (proměnné, které se objevují v předpokladech a důsledcích)
\item předpoklady - musí platit v daném stavu, aby se operátor dal použít
\item důsledky - literály (fluenty!), které budou platit, jakmile se operátor provede
\end{itemize}
\item akce je plně zinstanciovaný operátor (za proměnné jsme dosadili konstanty)
\item akce je relevantní pro cíl $g\equiv$ přispívá cíli $g$ a její důsledky (efekty) nejsou v~konfliktu s $g$
\item dále definujeme výsledek aplikace akce na stav a regresní množinu pro cíl a akci
\end{itemize}
\item terminologie
\begin{itemize}
\item popisu operátorů se obvykle říká doménový model
\item plánovací problém je trojice $(O,s_0,g)$, kde $O$ je doménový model, $s_0$ je výchozí stav a $g$ je cílová podmínka
\item plán je posloupnost akcí
\end{itemize}
\item plánování je realizováno prohledáváním stavového prostoru
\begin{itemize}
\item ten je mnohdy dost velký
\end{itemize}
\end{itemize}
\item dopředné a zpětné plánování
\begin{itemize}
\item forward (progression) planning
\begin{itemize}
\item postupujeme od výchozího k cílovému stavu
\item prohledávací algoritmus potřebuje dobrou heuristiku
\end{itemize}
\item backward (regression) planning
\begin{itemize}
\item prozkoumávají se jen relevantní akce $\to$ menší větvení než u forward planningu
\item používá množiny stavů spíše než individuální stavy $\to$ je těžké najít dobrou heuristiku
\end{itemize}
\end{itemize}
\end{itemize}
\subsection{Markovské rozhodovací procesy (MDP)}
\begin{itemize}
\item formulace problému (výpočet užitku, strategie)
\begin{itemize}
\item řešíme sekvenční rozhodovací problém
\item pro plně pozorovatelné stochastické (nedeterministické) prostředí
\item máme Markovův přechodový model a reward
\begin{itemize}
\item přechodový model $P(s'\mid s,a)$ \dots{} pravděpodobnost dosažení stavu $s'$, když ve stavu $s$ použiju akci $a\in A(s)$
\item odměna (reward) $R(s)$ \dots{} získá ji agent ve stavu $s$, je krátkodobá
\item utility function $U([s_0,s_1,s_2,\dots])=R(s_0)+\gamma R(s_1)+\gamma^2 R(s_2)+\dots$
\begin{itemize}
\item kde $\gamma$ je \textit{discount factor}, číslo mezi 0 a 1
\begin{itemize}
\item jak moc si ceníme budoucích odměn
\end{itemize}
\item utility je \quotedblbase{}dlouhodobá lokální odměna\textquotedblleft{}
\item hodnota utility funkce je konečná i pro nekonečnou posloupnost stavů, protože zjevně $U([s_0,\dots])\leq \frac{R_\text{max}}{1-\gamma}$
\end{itemize}
\end{itemize}
\item řešením MDP je policy (strategie) - funkce doporučující akci pro každý stav
\begin{itemize}
\item tzn. strategie je zobrazení $\pi:S\to A$
\item potřebujeme zobrazení, protože kdyby řešením byla fixní sekvence akcí, tak by to nefungovalo pro stochastická prostředí (mohlo by se stát, že skončíme v jiném stavu, než jsme mysleli)
\item optimální strategie $\equiv$ strategie, která vrací největší střední hodnotu užitku (utility)
\item optimální strategie nezávisí na počátečním stavu - je jedno, kde jsme začali, pro nalezení nejvhodnější další akce by měl být důležitý jen aktuální stav
\end{itemize}
\end{itemize}
\item Bellmanova rovnice
\begin{itemize}
\item opravdový užitek stavu je reward (odměna) stavu ve spojení se střední hodnotou užitku následného stavu $\to$ to vede na Bellmanovu rovnici
\item vezmu reward a discountovaný užitek okolí
\item akorát nevím, kterou akci provedu, tak vezmu všechny a maximum přes ně (násobím jejich užitek pravděpodobností)
\item $U(s)=R(s)+\gamma\max_a\sum_{s'}P(s'\mid s,a)\cdot U(s')$
\end{itemize}
\item iterace hodnot, iterace strategií
\begin{itemize}
\item předpoklad: známe reward $R$ a transition model $P$
\begin{itemize}
\item chceme najít užitek $U$ nebo optimální strategii $\pi$
\end{itemize}
\item soustava Bellmanových rovnic není lineární - obsahuje maximum
\begin{itemize}
\item můžeme je řešit aproximací
\item použijeme iterativní přístup
\begin{itemize}
\item nastavíme nějak užitky - třeba jako nuly
\item provedeme update - aplikujeme Bellmanovu rovnici
\item iterujeme, dokud to nezkonverguje
\end{itemize}
\item $\to$ \textbf{value iteration}
\end{itemize}
\item strategie se ustálí dřív než užitky
\begin{itemize}
\item policy loss \dots{} vzdálenost mezi optimálním užitkem a užitkem strategie
\item můžeme iterativně zlepšovat policy, dokud to jde
\item $\to$ \textbf{policy iteration}
\item z rovnic nám zmizí maximum $\to$ máme lineární rovnice
\begin{itemize}
\item policy evaluation = nalezení řešení těchto rovnic (tak najdeme utilities stavů)
\item gaussovka je $O(n^3)$
\item u velkých stavových prostorů dává smysl k policy evaluation použít value iteration
\end{itemize}
\item algoritmus doběhne, protože policies je jenom konečně mnoho a vždycky hledáme tu lepší
\end{itemize}
\end{itemize}
\item POMDP (základní definice)
\begin{itemize}
\item máme přechodový model $P(s'\mid s,a)$, akce $A(s)$ i odměny $R(s)$
\item navíc nám přibude sensor model $P(s\mid e)$
\item místo reálných stavů můžeme používat ty domnělé (belief states)
\item modely přechodů a senzorů jsou reprezentovány dynamickou bayesovskou sítí
\item přidáme rozhodování a užitky, čímž dostaneme dynamickou rozhodovací síť
\item generujeme si stromeček, kde se střídají vrstvy akcí a belief states
\item používá se algoritmus podobný algoritmu \textit{expected minimax}
\end{itemize}
\end{itemize}
\subsection{Hry a teorie her}
\begin{itemize}
\item algoritmy Minimax a alfa-beta prořezávání
\begin{itemize}
\item algoritmus minimax umožňuje u hry dvou hráčů s nulovým součtem a s úplnou informací určit optimální strategie obou hráčů
\begin{itemize}
\item procházíme strom hry
\item v listech nebo v určité hloubce získáme ohodnocení vrcholů
\begin{itemize}
\item omezení hloubky je důležité zvláště u her, jejichž stromy se hodně větví - je potřeba zvolit vhodný přístup k ohodnocení konkrétní herní situace
\end{itemize}
\item následně vždy v rodiči zvolíme minimum nebo maximum ze synů (v závislosti na tom, který hráč je zrovna na tahu - volba minima a maxima se střídají stejně jako tahy hráčů)
\item takto pro každý vrchol určíme, jaký tah je v dané situaci optimální
\begin{itemize}
\item za předpokladu, že protihráč hraje optimálně - což je nejhorší možný scénář
\end{itemize}
\end{itemize}
\item alfa-beta prořezávání
\begin{itemize}
\item je to rozšíření minimaxu, které může zrychlit jeho běh tak, že omezuje počet větví, které je potřeba projít a vyhodnotit
\item vycházíme z následujícího pozorování
\begin{itemize}
\item v algoritmu typicky hledáme maximum z několika hodnot, které jsou minimy (nebo naopak)
\item v určitém vrcholu máme seznam doposud určených minim z jednotlivých větví (z nich je nějaké maximální)
\item snažíme se určit další minimum, na základě jedné z podvětví máme nějaké prozatímní minimum (museli bychom projít další podvětve pod daným vrcholem, abychom dostali výsledné minimum)
\item pokud je tohle prozatímní minimum menší než prozatímní maximum (ve vyšší úrovni), tak už ty další větve nemusíme procházet, protože výsledné minimum nemůže být větší než prozatímní minimum
\end{itemize}
\item alfa-beta prořezávání nemá vliv na výsledek algoritmu
\item záleží na pořadí, ve kterém tahy (větve) procházíme
\item minimax se obvykle používá rekurzivně, meze určující prořezávání (pro jednoho hráče je to maximum, pro druhého minimum) se do funkce předávají jako parametry alfa a beta
\end{itemize}
\end{itemize}
\item základy teorie her (vězňovo dilema, Nashovo ekvilibrium)
\begin{itemize}
\item jednou z úloh teorie her je volba vhodné strategie - tedy návodu, jak hru hrát
\begin{itemize}
\item uvažujeme hry v normálním tvaru: hru tvoří hráči, jejich akce a užitkové funkce
\item hráči táhnou zároveň, každý provede jedinou akci
\item strategie může být čistá (přímo konkrétní akce) nebo smíšená (pravděpodobnostní rozdělení přes akce)
\end{itemize}
\item díváme se na optimalitu strategie nebo strategického profilu (množiny strategií všech hráčů) z různých úhlů pohledu
\begin{itemize}
\item strategie $s_1$ dominuje strategii $s_2$, pokud při volbě $s_1$ získá hráč větší užitek než při volbě $s_2$ nehledě na strategie ostatních hráčů
\begin{itemize}
\item strategie je dominantní, pokud dominuje všechny ostatní strategie
\end{itemize}
\item \textit{Nashovo ekvilibrium} je takový strategický profil, že by žádný hráč neměnil svou strategii, kdyby znal strategie ostatních
\item strategický profil $s$ \textit{Pareto dominuje} profil $s'$, pokud si změnou z $s'$ na $s$ žádný hráč nepohorší, ale aspoň jeden hráč si polepší
\begin{itemize}
\item strategický profil je \textit{Pareto optimální}, pokud neexistuje profil, který by ho Pareto dominoval
\end{itemize}
\item další zajímavým typem ekvilibria je korelované ekvilibrium
\end{itemize}
\item vězňovo dilema
\begin{itemize}
\item dva hráči (vězni)
\item vězeň může vypovídat (testify/defect) nebo odmítnout vypovídat (refuse/cooperate)
\begin{itemize}
\item oba vypovídají $\to$ oba ztrácejí málo (užitek -5)
\item jeden vypovídá, druhý ne $\to$ jeden neztratí nic (užitek 0), druhý hodně ztratí (užitek -10)
\item oba odmítnou $\to$ oba ztrácejí velmi málo (užitek -1)
\end{itemize}
\item \textit{vypovídat} je čistá dominantní strategie
\begin{itemize}
\item ale když vypovídají oba hráči, tak vyhrála policie - lepší by bylo, kdyby oba odmítli vypovídat
\item našli jsme Nashovo ekvilibrium - nikdo neprofituje ze změny strategie, pokud druhý hráč zůstane u stejné strategie
\end{itemize}
\end{itemize}
\end{itemize}
\item mechanism design (typy aukcí)
\begin{itemize}
\item jak se pozná dobrý mechanismus aukce
\begin{itemize}
\item maximalizuje výnos pro prodejce
\item vítěz aukce je agent, který si věci nejvíc cení
\item zájemci by měli mít dominantní strategii
\end{itemize}
\item anglická aukce (ascending-bid)
\begin{itemize}
\item začnu s $b_\text{min}$, pokud je to nějaký zájemce ochotný zaplatit, tak se ptám na $b_\text{min}+d$
\item strategie: přihazuju, dokud cena není vyšší než moje hodnota
\begin{itemize}
\item dominantní strategie
\end{itemize}
\end{itemize}
\item holandská aukce
\begin{itemize}
\item začíná se vyšší cenou, snižuje se, dokud ji někdo nepřijme
\end{itemize}
\item obálková aukce
\begin{itemize}
\item největší nabídka vítězí
\item neexistuje dominantní strategie
\begin{itemize}
\item pokud věříš, že maximum ostatních nabídek je $b_0$, tak bys měl nabídnout $b_0+\varepsilon$, pokud je to menší částka než hodnota, kterou pro tebe věc má
\end{itemize}
\end{itemize}
\item obálková second-price aukce (Vickrey)
\begin{itemize}
\item vyhraje ten, kdo nabídne nejvyšší cenu, platí druhou nejvyšší nabídnutou cenu
\item dominantní strategie je tam dát svoji hodnotu
\end{itemize}
\item VCG mechanismus
\begin{itemize}
\item lze použít při obecnějších (víceparametrických) aukcích
\end{itemize}
\end{itemize}
\end{itemize}
\subsection{Strojové učení}
\begin{itemize}
\item základní druhy učení (s učitelem, bez učitele, zpětnovazební)
\begin{itemize}
\item učení bez učitele (unsupervised learning) - agent se učí vzory ve vstupu, aniž by měl nějakou explicitní zpětnou vazbu
\item zpětnovazební učení (reinforcement learning) - agent dostává odměny nebo tresty podle toho, jak se chová v prostředí
\item učení s učitelem (supervised learning) - agent se učí funkci, která mapuje vstupy na výstupy
\end{itemize}
\item rozhodovací stromy (definice, konstrukce)
\begin{itemize}
\item přijímá vektor hodnot atributů, vrací výslednou hodnotu
\item na základě tabulky příkladů se dá postavit strom
\begin{itemize}
\item každý vnitřní vrchol odpovídá testu jednoho atributu
\item dá se zjistit odůvodnění konkrétního rozhodnutí
\item strom se dá postavit hladovou metodou rozděl a panuj
\begin{itemize}
\item k dělení se používá entropie
\end{itemize}
\item každý list určuje hodnotu, kterou klasifikační funkce vrátí
\end{itemize}
\end{itemize}
\item regrese, SVM (základní principy)
\begin{itemize}
\item při regresi přiřazujeme položce $x$ vhodnou funkční hodnotu $y$
\begin{itemize}
\item je to jedna z úloh strojového učení s učitelem
\end{itemize}
\item support-vector machine (SVM)
\begin{itemize}
\item stojí na lineární regresi
\item pokud lze třídy oddělit nadrovinou, zvolí takovou, která je nejdál od všech dat (příkladů)
\item nadrovina \dots{} maximum margin separator
\item pokud nejde použít nadrovinu, provede mapování do vícedimenzionálního prostoru (pomocí kernel function), kde už příklady půjde oddělit
\item SVMs jsou neparametrická metoda - příklady blíže k separátoru jsou důležitější, říká se jim support vectors (vlastně definují ten separátor)
\end{itemize}
\end{itemize}
\item Bayesovské učení, EM algoritmus
\begin{itemize}
\item bayesovské učení
\begin{itemize}
\item máme několik hypotéz
\item na základě pozorování upravujeme jejich pravděpodobnost
\item \textit{maximálně věrohodná hypotéza}
\begin{itemize}
\item z množiny hypotéz je to ta, která přiřazuje naměřeným datům největší pravděpodobnost
\item $\mathrm{argmax}_h\,p(\mathrm{data}\mid h)$ nebo také $\mathrm{argmax}_h\,p_h(\mathrm{data})$
\item pokud je hypotéza parametrizovaná (hledáme parametr), tak můžeme hledat hodnotu parametru, pro níž je derivace logaritmu pravděpodobnosti nulová
\begin{itemize}
\item tomu se říká maximum-likelihood parameter learning
\end{itemize}
\end{itemize}
\item z $p(\mathrm{data}\mid h)$ můžeme získat $p(h\mid\mathrm{data})=\frac{p(\mathrm{data}\mid h)p(h)}{\sum_{h'} p(\mathrm{data}\mid h')p(h')}$
\begin{itemize}
\item pokud mají všechny hypotézy stejné apriorní pravděpodobnosti, stačí vzít $p(\mathrm{data}\mid h)$ a normalizovat je, jelikož $\frac{p(h)}{\sum_{h'} p(\mathrm{data}\mid h')p(h')}$ bude konstantní
\end{itemize}
\item \textit{bayesovsky optimální predikce} pravděpodobnosti pro vstupní data a konkrétní predikci $z$
\begin{itemize}
\item $p(z\mid\mathrm{data})=\sum_h p(z\mid h)\cdot p(h\mid\mathrm{data})$
\end{itemize}
\end{itemize}
\item algoritmus expectation-maximization (EM)
\begin{itemize}
\item používá se k maximum-likelihood odhadu, pokud se nedá jednoduše spočítat
\item předstíráme, že známe parametry modelu (na začátku je určíme náhodně)
\item iterujeme dva kroky (dokud to nezkonverguje)
\begin{itemize}
\item dopočteme střední hodnoty skrytých proměnných (E-step, expectation)
\item upravíme parametry, abychom maximalizovali likelihood modelu (M-step, maximization)
\end{itemize}
\item příklad
\begin{itemize}
\item chceme odhadnout průměrnou výšku mužů a žen
\item máme data o naměřených výškách, ale chybí nám údaje o pohlaví (zda je to výška muže nebo ženy)
\begin{itemize}
\item třídy označíme jako 0 (muži), 1 (ženy)
\end{itemize}
\item začneme tak, že oba parametry (průměrná výška mužů $\mu_0$, průměrná výška žen $\mu_1$) inicializujeme náhodně
\item v E-kroku spočítáme střední hodnoty skrytých proměnných
\begin{itemize}
\item pohlaví je skrytá proměnná (pro každou naměřenou výšku)
\item střední hodnota bude odpovídat pravděpodobnosti, že konkrétní výška patří ženě
\end{itemize}
\item v M-kroku upravíme hodnoty $\mu_0,\mu_1$, abychom maximalizovali likelihood modelu
\begin{itemize}
\item tedy spočítáme vážený průměr výšek
\item váhy odpovídají pravděpodobnosti, že výška patří muži (u $\mu_0$) / ženě (u $\mu_1$)
\end{itemize}
\item poznámka: může se jednoduše stát, že nám parametry $\mu_0,\mu_1$ vyjdou naopak (EM je metoda učení bez učitele, hledá vzory v datech, nezajišťuje jejich intepretaci)
\end{itemize}
\end{itemize}
\end{itemize}
\item pasivní zpětnovazební učení (definice, metody ADP a TD)
\begin{itemize}
\item uvažujeme agenta, který se účastní MDP (markovského rozhodovacího procesu)
\item pasivní učení - známe strategii $\pi$ (je pevně daná) a učíme se, jak je dobrá (učíme se utility funkci)
\item agent nezná přechodový model ani reward function
\item přímý odhad utility
\begin{itemize}
\item máme trace (běh) mezi stavy, z toho počítáme utility function
\item odhadujeme užitky jednotlivých stavů
\begin{itemize}
\item užitek stavu = \textit{reward-to-go} (tedy odměna tohoto a všech budoucích stavů)
\item odměny z budoucích stavů bychom mohli zohleďnovat s discount factorem $\gamma$ (viz výše)
\end{itemize}
\item užitky pro jednotlivé stavy prostě průměrujeme
\item v podstatě učení s učitelem - na vstupu je stav, na výstupu užitek stavu (neboli \textit{reward-to-go})
\item nevýhoda
\begin{itemize}
\item utilities nejsou nezávislé, ale řídí se Bellmanovými rovnicemi
\item hledáme v prostoru hypotéz, který je výrazně větší, než by musel být (obsahuje spoustu funkcí, co porušují Bellmanovy rovnice) $\to$ algoritmus konverguje pomalu
\end{itemize}
\end{itemize}
\item adaptivní dynamické programování (ADP)
\begin{itemize}
\item agent se učí přechodový model a odměny
\item přechodový model odhaduje přímo z pozorovaných přechodů
\item utilitu počítá z Bellmanových rovnic třeba pomocí value iteration
\begin{itemize}
\item využíváme toho, že užitky nejsou nezávislé
\end{itemize}
\item ADP zajišťuje, že jsou odhady užitků konzistentní
\end{itemize}
\item temporal-difference (TD) learning
\begin{itemize}
\item upravuje stav, aby odpovídal aktuálně pozorovanému následníkovi
\item když dojde k přechodu ze stavu $s$ do stavu $s'$, tak na $U^\pi(s)$ použijeme update: $U^\pi(s)\leftarrow U^\pi(s)+\alpha\cdot (R(s)+\gamma\cdot U^\pi(s')-U^\pi(s))$
\begin{itemize}
\item $\alpha$ je learning rate
\end{itemize}
\item vlastnosti
\begin{itemize}
\item nepotřebuju přechodový model (je to model-free metoda)
\item konverguje to pomaleji než ADP
\end{itemize}
\end{itemize}
\end{itemize}
\item aktivní zpětnovazební učení (definice, explorace vs. exploitace, Q-učení, SARSA)
\begin{itemize}
\item uvažujeme agenta, který se účastní MDP (markovského rozhodovacího procesu)
\item aktivní učení - agent se učí strategii (policy; jak se rozhodovat) a taky utility funkci
\item mohli bychom použít ADP, tím vznikne hladový agent, který opakuje zažité vzory (nezkouší nové věci)
\item jak může volba optimální akce vést k suboptimálním výsledkům?
\begin{itemize}
\item naučený model není stejný jako reálné prostředí
\item akce nejsou pouze zdrojem odměn, ale také přispívají k učení - ovlivňují vstupy
\item zlepšováním modelu se postupně zvětšují odměny, které agent získává
\end{itemize}
\item je potřeba najít optimum mezi exploration a exploitation
\begin{itemize}
\item pure exploration - agent nepoužívá naučené znalosti
\item pure exploitation - riskujeme, že bude agent pořád dokola opakovat zažité vzory
\item základní myšlenka: na začátku preferujeme exploration, později lépe rozumíme světu, takže nepotřebujeme tolik prozkoumávat
\end{itemize}
\item exploration policies
\begin{itemize}
\item první možná exploration policy: agent si zvolí náhodnou akci s pravděpodobností $\frac1t$ (kde $t$ je čas), jinak se řídí hladovou strategií
\begin{itemize}
\item nakonec to konverguje k optimální strategii, ale může to být extrémně pomalé
\end{itemize}
\item rozumnější přístup je přiřadit váhu akcím, které agent ještě nevyzkoušel
\item existuje alternativní TD metoda, říká se jí Q-learning
\begin{itemize}
\item $Q(s,a)$ označuje hodnotu toho, že provedeme akci $a$ ve stavu $s$
\item q-hodnoty jsou s utilitou ve vztahu $U(s)=\max_a Q(s,a)$
\item můžeme napsat omezující podmínku $Q(s,a)=R(s)+\gamma\sum_{s'} P(s'\mid s,a)\cdot \max_{a'}Q(s',a')$
\begin{itemize}
\item tohle vyžaduje, aby se model naučil i $P(s'\mid s,a)$
\end{itemize}
\item Q-learning ale nevyžaduje model přechodů - je to bezmodelová (model-free) metoda, potřebuje akorát q-hodnoty
\item $Q(s,a)\leftarrow Q(s,a)+\alpha\cdot (R(s)+\gamma\cdot\max_{a'}Q(s',a')-Q(s,a))$
\begin{itemize}
\item počítá se, když je akce $a$ vykonána ve stavu $s$ a vede do stavu $s'$
\end{itemize}
\end{itemize}
\item state-action-reward-state-action (SARSA)
\begin{itemize}
\item je to varianta Q-learningu
\item $Q(s,a)\leftarrow Q(s,a)+\alpha\cdot (R(s)+\gamma\cdot Q(s',a')-Q(s,a))$
\begin{itemize}
\item pravidlo se aplikuje na konci pětice $s,a,r,s',a'$, tedy po aplikaci akce $a'$
\end{itemize}
\end{itemize}
\item pro hladového agenta (který volí podle největšího $Q$) jsou algoritmy SARSA a základní Q-learning stejné
\begin{itemize}
\item rozdíl je vidět až u agenta, který není úplně hladový, ale provádí i průzkum (exploration)
\end{itemize}
\item u SARSA se bere v úvahu reálně zvolená akce
\begin{itemize}
\item SARSA je on-policy - jeho strategie se upravuje přímo za běhu algoritmu
\item Q-learning je off-policy - algoritmus se učí optimální strategii, ale během učení se může používat úplně jiná strategie
\end{itemize}
\end{itemize}
\item při aktivním učení je výhodnější učit se q-hodnoty než užitkovou funkci
\begin{itemize}
\item abychom našli nejvýhodnější akci v daném stavu pomocí utility, museli bychom při výpočtu použít pravděpodobnost, že danou akcí přejdeme do určitého stavu
\item q-hodnota nám dá přímo řekne, jak je užitečné provést danou akci - k volbě nejlepší akce nepotřebujeme přechodový model
\end{itemize}
\end{itemize}
\end{itemize}
